\documentclass[11pt]{article}
\usepackage[a4paper,margin=1in]{geometry}
\usepackage{fontspec}
\usepackage[boldfont=false]{xeCJK}
\setCJKmainfont{FandolSong-Regular.otf}
\usepackage{amsmath}
\usepackage{amssymb}
\usepackage{array}
\usepackage{booktabs}
\usepackage{graphicx}
\usepackage{adjustbox}
\usepackage{enumitem}
\setlist[itemize]{topsep=2pt,partopsep=0pt,parsep=0pt,itemsep=2pt,leftmargin=1.4em}
\setlist[enumerate]{topsep=2pt,partopsep=0pt,parsep=0pt,itemsep=2pt,leftmargin=1.6em}
\usepackage{parskip}
\usepackage{fvextra}
\fvset{breaklines=true,breakanywhere=true,fontsize=\small}
\setlength{\parindent}{0pt}

\begin{document}

\begin{center}
  {\LARGE\bfseries Control and coordination}\\[0.35em]
  {\large A Level Biology}
\end{center}
\vspace{0.6em}

\section*{Two systems for control}

The body has two coordination systems. The \textbf{endocrine system} 内分泌系统 sends chemical \textbf{hormones} 激素 (such as ADH, glucagon and insulin) in the blood. The \textbf{nervous system} 神经系统 sends fast electrical signals along nerve cells.

\begin{center}
\adjustbox{max width=\linewidth}{%
\begin{tabular}{lll}
\toprule
\textbf{Feature} & \textbf{Nervous system} & \textbf{Endocrine system} \\
\midrule
signal & electrical \textbf{impulse} 冲动 & chemical hormone \\
transport & along nerve cells & in the blood \\
speed & very fast & slower \\
how long it lasts & short & longer \\
\bottomrule
\end{tabular}}
\end{center}

\section*{Neurones}

A \textbf{neurone} 神经元 is a nerve cell. There are three kinds:

\begin{itemize}
\item a \textbf{sensory neurone} 感觉神经元 carries impulses from a receptor towards the brain or spinal cord.

\item a \textbf{motor neurone} 运动神经元 carries impulses out to an effector, such as a \textbf{muscle} 肌肉.

\item \textbf{intermediate neurones} 中间神经元 connect sensory neurones to motor neurones.

\end{itemize}

A neurone has a long fibre (the axon) along which the impulse travels.

\section*{Detecting a stimulus}

A \textbf{sensory receptor} 感受器 cell detects a \textbf{stimulus} 刺激 and starts an impulse in a sensory neurone. For example, a \textbf{chemoreceptor} 化学感受器 cell in a \textbf{taste bud} 味蕾 detects chemicals in food, and this triggers an \textbf{action potential} 动作电位 in the sensory neurone.

\section*{The nerve impulse}

The impulse is a change in the voltage across the neurone's membrane.

\begin{itemize}
\item \textbf{resting potential} 静息电位 — when no impulse passes, the inside is negative compared to the outside. This is maintained by a pump that moves sodium ions out and potassium ions in, and by the membrane being more permeable to potassium.

\item \textbf{action potential} — a stimulus makes sodium channels open, so sodium ions rush in and the inside briefly becomes positive (depolarisation). Then potassium channels open, potassium ions leave, and the \textbf{membrane potential} 膜电位 returns to negative (repolarisation).

\item \textbf{refractory period} 不应期 — just after an action potential, the sodium channels cannot open again for a short time. This restores the resting potential and sets a limit on how often impulses can be sent (their frequency).

\end{itemize}

\subsection*{Faster impulses: myelin}

Some neurones are wrapped in a fatty \textbf{myelin sheath} 髓鞘. The impulse cannot cross the sheath, so it jumps from one gap to the next. This jumping, called \textbf{saltatory conduction} 跳跃式传导, makes the impulse travel much faster.

\section*{The synapse}

A \textbf{synapse} 突触 is a tiny gap between two neurones. A \textbf{cholinergic synapse} 胆碱能突触 passes the signal like this:

\begin{enumerate}
\item the impulse arrives and makes \textbf{calcium ions} 钙离子 enter the first neurone.

\item this makes vesicles release a \textbf{neurotransmitter} 神经递质 called \textbf{acetylcholine} 乙酰胆碱.

\item the acetylcholine diffuses across the gap and binds to receptors on the next neurone.

\item this starts a new impulse in the next neurone.

\end{enumerate}

\section*{Muscles and how they contract}

A \textbf{neuromuscular junction} 神经肌肉接头 is like a synapse, but between a motor neurone and a muscle.

\textbf{Striated muscle} 横纹肌 is made of long fibres. Each fibre is divided into repeating units called \textbf{sarcomeres} 肌节. Inside are two kinds of \textbf{filament} 肌丝: thick ones (myosin) and thin ones (actin). The \textbf{T-tubule} T小管 system carries the impulse deep into the fibre, and the \textbf{sarcoplasmic reticulum} 肌质网 stores and releases calcium ions.

The \textbf{sliding filament model} 肌丝滑动模型 explains contraction:

\begin{enumerate}
\item an impulse causes the sarcoplasmic reticulum to release calcium ions.

\item the calcium ions bind to \textbf{troponin} 肌钙蛋白, which makes \textbf{tropomyosin} 原肌球蛋白 move and uncover the binding sites on the actin.

\item the myosin heads attach to the actin and pull it inwards, using energy from ATP.

\item the thin filaments slide over the thick ones, so each sarcomere gets shorter and the muscle contracts.

\end{enumerate}

\section*{Control and coordination in plants}

\subsection*{The Venus fly trap}

The \textbf{Venus fly trap} 捕蝇草 has tiny hairs on its trap-like leaves. When an insect touches the hairs, a fast electrical signal spreads, the cells quickly lose water and change shape, and the trap snaps shut to catch the insect.

\subsection*{Auxin and growth}

\textbf{Auxin} 生长素 is a plant hormone that makes cells grow longer (\textbf{elongation} 伸长 growth). It does this by making the cell pump \textbf{protons} 质子 (hydrogen ions) into the \textbf{cell wall} 细胞壁. This \textbf{acidifies} 酸化 the wall, which loosens it, so the cell can stretch as water enters.

\subsection*{Gibberellin and germination}

\textbf{Gibberellin} 赤霉素 controls the \textbf{germination} 萌发 of \textbf{barley} 大麦 seeds. It switches on the genes that make amylase, which then breaks down stored starch into sugars to feed the growing seedling.


\end{document}
